\section{A minimal-character scale obstruction in the marked $j=0$ packet}

The cyclic cubic strategy admits a particularly sharp exact test in the
marked $j=0$ family. Put
\[
q(t)=t^2-t+1,\qquad r(t)=t(t-1),
\]
and specialize the marking parameter to $\mu=2$, so that
\[
a=q(2)=3,\qquad b=\frac{q(2)^3}{r(2)}=\frac{27}{2}.
\]
After the cyclic base change
\[
u^3=c\,t(t-1),\qquad c\in\Q^*,
\]
the surface over $K=\Q(\sqrt{-3})$ is
\[
y^2=x^3+\frac{b^2}{c^2}u^6-a^3\left(\frac{u^3}{c}+1\right)^3.
\]

\begin{theorem}[Minimal-character scale obstruction]
A character-one section in the minimal polynomial box
\[
x=u(Au^3+B),\qquad y=Cu^6+Du^3+E,
\qquad A,B,C,D,E\in K,
\]
can exist only when
\[
[c]\in\{[1],[2]\}\subset\Q^*/\Q^{*3}.
\]
Both classes occur and carry a non-torsion Eisenstein eigenline over $K$.
However, the corresponding genus-one bases have
\[
\rank C_1(\Q)=\rank C_2(\Q)=0.
\]
Thus no positive-rank cyclic base in this $\mu=2$ family carries a minimal
character-one eigensection.
\end{theorem}

\begin{proof}
Coefficient comparison gives
\[
\begin{aligned}
C^2&=A^3,&
2CD&=3A^2B-\frac{27}{c^3},\\
D^2+2CE&=3AB^2+\frac{405}{4c^2},&
2DE&=B^3-\frac{81}{c},\\
E^2&=-27.&&
\end{aligned}
\]
The branch $B=0$ forces
$D=3E/(2c)$ and $C=E/(3c^2)$, whence
$D^2+2CE=-315/(4c^2)$, a contradiction. Hence $B\ne0$.
Choose $E=3\sqrt{-3}$ and set
\[
\rho=\frac{cA}{B},\qquad S=cB^3,
\qquad g=\frac{c^2C}{\sqrt{-3}},
\qquad d=\frac{cD}{\sqrt{-3}}.
\]
Exact lexicographic elimination gives
\[
(2\rho+1)^3(10\rho^3-48\rho^2+12\rho-1)=0.
\]
The cubic factor is irreducible over $\Q$, so it has no root in the quadratic
field $K$. Therefore $\rho=-1/2$. The two remaining normalized solutions
are
\[
(S,g,d)=\left(54,-\frac32,\frac32\right),
\qquad
(S,g,d)=\left(216,3,-\frac{15}{2}\right).
\]
If $z\in K$ and $z^3\in\Q$, then $z^3\in\Q^3$; hence
$K^3\cap\Q=\Q^3$. Since $S/c=B^3$, the class of $c$ equals the class of
$S$. Now $54=2\cdot3^3$ and $216=6^3$, proving the two-class statement.
Exact substitution supplies sections for both classes. Their specializations
at $u=1$ have reduction orders $(13,19)$ at primes above $(7,13)$ in class
$[1]$, and $(13,21)$ in class $[2]$. The good-reduction constraints on a
hypothetical torsion order are incompatible, so the sections are non-torsion.
Finally, SageMath with proof flags computes exact rank zero for the two base
curves, which are birational to $Y^2=X^3+16c^4$.
\end{proof}

This theorem explains an earlier computational split. The scale $c=27$ is
cube-equivalent to $c=1$ and exhibits minimal eigensections, but its base has
rank zero. The scale $c=36$ has base rank one, but the complete minimal
scheme is empty. The incompatibility is structural for this marked surface,
not a failure of a small coefficient box.

\subsection{Decisive change of setting}

The minimal ansatz is exhausted. The next exact experiment is a controlled-
pole character-one search on the rank-one base $C_3$, followed in parallel by
the same Kummer classification along the CM-compatible marking conic
$q(\mu)=3s^2$. A successful one-pole orbit would contribute a genuine
Eisenstein line over a positive-rank genus-one base; failure would give a
second, sharper visibility obstruction.
